\chapter{表示学习、自编码器与重构误差}

\section{案例目标}

前面两章我们已经建立了两条重要直觉：一是神经网络可以表示明显非线性的映射，二是卷积结构适合图像和光谱中的局部模式。但现代 AI 里还有另一条同样重要的主线：\textbf{在没有明确类别标签的情况下，模型能不能先学到一个有用的内部表示？}

本章的目标有三点：

\begin{enumerate}
    \item 理解表示学习为什么不仅仅是“降维”，而是学习可复用的隐变量结构；
    \item 理解自编码器怎样通过重构任务学习正常数据的低维表示；
    \item 理解为什么重构误差可以成为一个最小异常分数，为后续更真实的光谱和图像工作流做准备。
\end{enumerate}

\section{为什么要从 PCA 走向自编码器}

在第 23 章里，我们已经用 PCA 看过一种经典的低维表示方法：把高维输入投影到少数几个主方向上，再用这些主方向解释数据的大尺度变化。PCA 的价值非常大，但它仍然有两个局限：

\begin{itemize}
    \item 它本质上是线性投影，难以表达更复杂的弯曲数据流形；
    \item 它更强调“方差最大的方向”，不一定直接对应最适合重构或下游任务的表示。
\end{itemize}

自编码器可以看作是从 PCA 迈向更灵活表示学习的一座桥。它不再直接求一个固定投影，而是通过一个\textbf{编码器}把输入压缩到低维瓶颈，再通过一个\textbf{解码器}把瓶颈表示还原回原始输入。

\section{一个极小的光谱重构教学数据集}

配套 notebook 使用 \nolinkurl{spectral_autoencoder_demo.csv}。它是一个教学版的一维光谱窗口数据集，每个样本包含 12 个连续 flux bin。数据表包括：

\begin{itemize}
    \item \texttt{sample\_id}；
    \item \texttt{split}；
    \item \texttt{spectrum\_role}；
    \item \texttt{pattern\_label}；
    \item \texttt{continuum\_slope} 与 \texttt{balmer\_depth} 两个教学用物理代理量；
    \item \texttt{f00} 到 \texttt{f11} 的归一化 flux 采样值。
\end{itemize}

这个数据集刻意构造成两部分：

\begin{itemize}
    \item 训练集只包含“正常”光谱窗口，它们主要由连续谱斜率和 Balmer 吸收深度这两个低维因素控制；
    \item 测试集同时包含正常样本和异常样本，异常模式包括局部发射突增、宇宙线尖峰、红端异常鼓包和过宽吸收结构。
\end{itemize}

这样设计的目的，是让学生直接看到：\textbf{如果正常数据真的主要落在一个低维流形上，那么一个小瓶颈自编码器就能学会重构它；而偏离这个流形的样本，往往会表现出更大的重构误差。}

\section{先做一个平均谱 baseline}

即使进入表示学习，本书仍然不跳过 baseline。notebook 的第一个对照组，是把训练集的平均谱当成“所有样本的统一重构模板”。最小形式可以写成：

\begin{examplebox}[平均谱 baseline]
\begin{lstlisting}[style=pythonstyle]
reconstruction = mean_train_spectrum
\end{lstlisting}
\end{examplebox}

这个 baseline 很粗糙，但它非常有教学价值：

\begin{itemize}
    \item 它提醒学生“重构”并不自动等于“学到了表示”；
    \item 它说明如果模型没有能力区分不同正常样本的内部结构，就只能输出一个模糊模板；
    \item 它给后面的自编码器提供了一个最小对照组。
\end{itemize}

\begin{table}[htbp]
\centering
\small
\resizebox{\textwidth}{!}{
\begin{tabular}{p{0.20\textwidth}ccp{0.32\textwidth}}
\toprule
方法 & 正常测试样本平均重构误差 & 异常测试样本平均重构误差 & 教学解读 \\
\midrule
平均谱 baseline & 0.85 & 4.97 & 只能记住一个粗模板，保不住正常变化 \\
二维瓶颈自编码器 & 0.01 & 5.25 & 学到正常流形，同时保留异常高误差 \\
共享局部 patch 自编码器 & 0.00 & 2.26 & 共享局部权重，可定位异常窗 \\
端到端 \texttt{Conv1d} 自编码器 & 0.07 & 3.38 & 用整谱训练共享卷积核，兼看局部与整体偏离 \\
\bottomrule
\end{tabular}
}
\caption{本章教学 notebook 中几种最小重构方案的对照。重点不是追求最低数字，而是观察模型是否真的学会了正常数据中的可复用结构。}
\label{tab:ch33_mean_vs_autoencoder}
\end{table}

\section{自编码器到底在做什么}

最小自编码器由两部分组成：

\begin{itemize}
    \item 编码器：把高维输入压缩到低维 latent 向量；
    \item 解码器：从 latent 向量重构原始输入。
\end{itemize}

在本章 notebook 中，我们使用的是一个最小二维瓶颈：

\begin{examplebox}[一个最小自编码器]
\begin{lstlisting}[style=pythonstyle]
latent = tanh(W_enc * x + b_enc)
reconstruction = W_dec * latent + b_dec
\end{lstlisting}
\end{examplebox}

这里最关键的教学点不是“层多不多”，而是\textbf{瓶颈约束}。当 latent 只有 2 个维度时，模型就被迫用非常紧凑的方式概括正常样本中最可重复的结构。

在当前教学数据上，notebook 还会进一步检查 latent 轴与两个教学代理量的关系：第一维主要沿着连续谱斜率展开，第二维则更多对应 Balmer 吸收深度的变化。这让学生更容易把“抽象表示”重新翻译回可解释的光谱变化方向。

\section{为什么这已经是一种自监督学习}

这个案例虽然很小，但它已经体现了自监督学习最核心的想法：\textbf{训练目标来自输入本身，而不是来自外部标签。}

在当前 notebook 中：

\begin{itemize}
    \item 训练时并不使用“正常”或“异常”的标签去做分类；
    \item 模型只被要求把输入光谱尽量重构回来；
    \item 是否为异常，是在训练完成后通过重构误差来判断的。
\end{itemize}

这与后面更现代的 masked modeling、对比学习和科学基础模型是同一条思想链上的不同复杂度版本。

\section{重构误差为什么能当异常分数}

如果训练集只包含正常样本，那么自编码器最擅长重构的，就应该是那些仍然落在“正常流形”附近的输入；相反，如果某个样本含有训练中没见过的尖峰、局部发射或过宽结构，它在低维瓶颈中就更难被忠实表示。

因此，本章 notebook 使用了一个最小异常分数：

\begin{examplebox}[最小异常分数]
\begin{lstlisting}[style=pythonstyle]
anomaly_score = mean((x - reconstruction) ** 2)
\end{lstlisting}
\end{examplebox}

这当然不是异常检测的终点，但它很好地说明了一个关键思想：\textbf{异常不一定要先有类别标签，很多时候可以先被定义为“偏离正常表示结构的样本”。}

notebook 的后半段还会把一部分正常样本临时拿出来充当 validation 集，用它校准“正常可接受误差”的范围，再把测试样本路由到 \texttt{normal\_like}、\texttt{manual\_review} 和 \texttt{anomaly\_candidate} 三个队列。在当前教学数据上，这个最小 workflow 会得到 3 个 \texttt{normal\_like}、2 个 \texttt{manual\_review} 和 3 个 \texttt{anomaly\_candidate}；其中一个需要人工复核的是正常样本在红端浅吸收边界上的延伸，另一个则是没有尖峰那么极端、但仍值得检查的宽吸收异常。

图~\ref{fig:ch33_normal_vs_anomaly} 给出了本章教学数据中的几类代表性模式。

\begin{figure}[htbp]
\centering
\begin{tikzpicture}
\begin{axis}[
    width=0.82\textwidth,
    height=0.54\textwidth,
    xlabel={flux bin},
    ylabel={normalized flux},
    xmin=0, xmax=11,
    ymin=0.68, ymax=1.48,
    xtick={0,2,4,6,8,10},
    grid=both,
    major grid style={draw=gray!20},
    minor grid style={draw=gray!10},
    tick label style={font=\small},
    label style={font=\small},
    legend style={font=\small, draw=none, fill=white, fill opacity=0.85, text opacity=1, at={(0.02,0.98)}, anchor=north west},
]
\addplot[mark=*, color=blue!70!black] coordinates {
    (0,1.085) (1,1.063) (2,1.026) (3,0.968) (4,0.867) (5,0.779)
    (6,0.763) (7,0.819) (8,0.888) (9,0.914) (10,0.919) (11,0.909)
};
\addlegendentry{正常：蓝端更亮，深吸收}

\addplot[mark=square*, color=orange!85!black] coordinates {
    (0,0.861) (1,0.881) (2,0.891) (3,0.887) (4,0.854) (5,0.832)
    (6,0.858) (7,0.932) (8,1.017) (9,1.073) (10,1.115) (11,1.147)
};
\addlegendentry{正常：红端更亮，浅吸收}

\addplot[mark=triangle*, color=red!70!black] coordinates {
    (0,0.979) (1,0.977) (2,0.964) (3,0.935) (4,0.873) (5,1.272)
    (6,1.126) (7,0.885) (8,0.955) (9,0.992) (10,1.013) (11,1.023)
};
\addlegendentry{异常：局部发射突增}

\addplot[mark=diamond*, color=purple!70!black] coordinates {
    (0,0.894) (1,0.906) (2,1.457) (3,0.890) (4,0.837) (5,0.796)
    (6,0.816) (7,0.897) (8,0.990) (9,1.047) (10,1.086) (11,1.114)
};
\addlegendentry{异常：宇宙线尖峰}
\end{axis}
\end{tikzpicture}
\caption{本章教学数据中的正常与异常光谱窗口示意。正常样本的主要变化可以被少数潜在因素概括，而异常样本会引入训练阶段未见过的局部结构。}
\label{fig:ch33_normal_vs_anomaly}
\end{figure}

\section{配套 notebook 做了什么}

本章 notebook 延续了前几章的克制风格：

\begin{itemize}
    \item 先读取小型光谱窗口数据，并区分训练集与测试集；
    \item 用平均谱做最小重构 baseline；
    \item 对输入做标准化；
    \item 用纯 Python 训练一个二维瓶颈的极小 \texttt{tanh} 自编码器；
    \item 比较正常样本和异常样本的重构误差；
    \item 再打印若干样本的 latent 坐标，检查它们是否与连续谱斜率、Balmer 深度等因素大致对应；
    \item 拿出一小部分正常样本做 validation 校准，再把测试样本路由到正常样本、人工复核和异常候选三个队列；
    \item 倒数第二段先用共享的 3-bin patch 自编码器滑过整段光谱，近似演示卷积自编码器的局部重构与 peak-patch 异常定位；
    \item 最后一段再把同样的共享权重思路推进到一个端到端训练的最小 \texttt{Conv1d} 自编码器；
    \item 然后把 dense bottleneck 的二维 latent 当成一个最小 normal reference library，做最近邻检索，并把 retrieved normal family 与 \texttt{Conv1d} residual window 合在一起形成 anomaly triage 摘要。
\end{itemize}

这种组织方式的好处是：

\begin{itemize}
    \item 学生可以看到表示学习并不一定要从很大很复杂的模型开始；
    \item “低维表示”和“异常分数”之间的关系被放在了同一个工作流里；
    \item validation 校准和 review queue 让学生看到：异常检测并不一定一步到位，更常见的是“先标出明显异常，再把边界样本交回人工复核”；
    \item 共享 patch 的扩展又把“局部感受野 + 权重共享 + 局部异常定位”顺手接回了第 32 章的一维卷积直觉；
    \item 端到端 \texttt{Conv1d} 自编码器又进一步说明：同样是共享局部结构，单独训练 patch 和整谱联合训练，看到的异常边界可以不同；
    \item 新增的 latent retrieval 段则说明：表示不仅可以用来打异常分数，也可以拿来做“最像哪类正常样本”的检索与 triage 上下文；
    \item notebook 同时把第 23 章的无监督学习、第 24 章的异常检测和第 32 章的光谱表示串了起来。
\end{itemize}

\section{再往前一步：共享局部编码器的卷积式扩展}

notebook 的最后一段没有直接训练完整的 \texttt{Conv1d} 自编码器，而是做了一个更容易讲清楚的近似版：先把每条光谱切成重叠的 3-bin patch，再用同一个小自编码器去重构所有 patch，最后把重叠区域平均回整条光谱。这样做已经把卷积表示里最关键的三件事接了起来：局部感受野、跨位置共享权重，以及由局部重构误差组成的全局异常分数。

在当前教学数据上，这个 shared-patch workflow 先用 workflow train 的正常样本训练局部编码器，再用 validation 正常样本同时校准整谱重构误差与峰值 patch 误差。校准后的阈值大约是：整谱 ready threshold 为 \texttt{0.0133}，peak-patch ready threshold 为 \texttt{0.0348}，更激进的 peak-patch anomaly threshold 为 \texttt{0.0811}。

这样路由后，测试集中的 3 个正常样本仍落在 \texttt{normal\_like}，红端浅吸收边界样本 \texttt{R16} 落在 \texttt{manual\_review}，而 4 个异常样本都进入 \texttt{anomaly\_candidate}。特别是原先在全局二维瓶颈 workflow 中仍停留在人工复核的 \texttt{A04}，在这一段里被提升为异常候选，因为它的宽吸收 trough 会在 \texttt{f07}--\texttt{f09} 一带连续触发较大的局部 patch mismatch。

这并不意味着 shared-patch 版本“绝对更好”，而是说明模型结构不同，看到的异常类型也会不同：全局瓶颈更像在学习正常流形，局部共享编码器则更像在扫描哪里出现了训练集中没见过的局部形状。

\section{再往前一步：一个端到端 Conv1d autoencoder}

notebook 的最后一小段再把 shared-patch 近似版往前推一步：这次不再先把 patch 各自训练好，而是直接用共享的 4-bin 卷积核在整条光谱上编码，再通过 overlap-add 的方式把所有局部表示拼回整谱重构。这样得到的是一个最小、但已经端到端训练的 \texttt{Conv1d} 自编码器。

在当前教学数据上，这个 end-to-end workflow 用 3 个共享 filter、\texttt{stride=1} 的局部编码器训练 workflow train 正常样本，再用 validation 正常样本校准整谱误差阈值。当前验证误差大约是 \texttt{0.0227}、\texttt{0.0666} 和 \texttt{0.1048}，据此得到 ready threshold \texttt{0.1363} 与 anomaly-candidate threshold \texttt{0.2620}。

按这个阈值路由后，测试集中的 3 个正常样本落在 \texttt{normal\_like}，\texttt{R16} 仍停在 \texttt{manual\_review}，而 4 个异常样本全部进入 \texttt{anomaly\_candidate}。其中最值得对照的是 \texttt{A04}：它在全局二维瓶颈 workflow 中仍是人工复核，在这里则被提升为异常候选，因为当共享卷积核被迫端到端重构整段宽吸收 trough 时，整谱误差会明显越过 validation 校准出来的异常阈值。

这个结果和上一段 shared-patch workflow 的教学意义并不重复。shared-patch 版本强调的是“局部感受野 + 权重共享 + patch mismatch”；而端到端 \texttt{Conv1d} 自编码器则进一步说明，局部模式一旦通过共享卷积核共同决定整谱重构，模型就会同时对局部线型和整体光谱一致性变得更敏感。

\section{再往前一步：latent retrieval 与 anomaly triage}

在补上 dense bottleneck、shared-patch 和端到端 \texttt{Conv1d} 三层工作流之后，notebook 又把最早那套二维 latent 表示重新拿回来，当成一个极小的 \textbf{normal reference library}。做法很简单：只把 workflow train 的正常样本放进 reference bank，然后对 validation 和 test 样本在 latent 空间里做最近邻检索。

这一段的重点不是拿 latent 最近邻直接替代异常检测，而是让学生看到两件互补的事：

\begin{itemize}
    \item latent retrieval 可以回答“这条候选光谱整体最像哪一类正常样本”；
    \item residual / reconstruction workflow 则回答“它具体在哪一段开始不像”。
\end{itemize}

在当前教学数据上，validation 正常样本的最近邻距离大约是 \texttt{0.2831}、\texttt{0.3440} 和 \texttt{0.4257}，据此可以校准出一个最小 latent support threshold \texttt{0.4683}。在这个阈值下，4 个正常测试样本都能检索回正确的正常 family，normal-test family accuracy 为 \texttt{1.00}。

但更有教学意义的是异常样本：4 个异常里有 3 个仍然落在这个 latent support 半径之内。也就是说，二维 bottleneck 仍然主要在编码连续谱斜率和 Balmer 深度这样的全局结构；像宇宙线尖峰、红端鼓包和宽吸收 trough 这种局部异常，并不会自动把样本从正常 latent manifold 上完全甩出去。

notebook 因而把 latent retrieval 和上一段 \texttt{Conv1d} residual window 合在一起，给出一个更接近真实科研分诊的最小 triage：

\begin{itemize}
    \item 3 个正常样本落在 \texttt{retrieval\_ready}；
    \item \texttt{R16} 落在 \texttt{normal\_manifold\_edge}，说明它整体仍像 \texttt{red\_continuum\_shallow\_line} 家族，只是在红端边界上超出了 validation 支撑；
    \item \texttt{A02}、\texttt{A03} 和 \texttt{A04} 虽然都还能检索回一个看似合理的正常 family，但因为局部 residual window 仍然明显异常，所以被标成 \texttt{localized\_anomaly\_over\_known\_family}；
    \item \texttt{A01} 则是更强的 \texttt{manifold\_shift\_anomaly\_candidate}，因为它不仅 residual 明显异常，连 latent 最近邻距离也已经越过 validation 校准出来的 support 半径。
\end{itemize}

这一步非常适合教学，因为它把“检索”和“异常检测”这两件常被混在一起的事情分开了：\textbf{retrieval 更像是在找上下文和参照物，而 anomaly gate 更像是在判断这种相似性是否还能被信任。}

\section{这和真实表示学习还有多远}

当然，本章的教学自编码器还远不是现实中的完整表示学习系统。和真实工作流相比，它至少还缺少几层内容：

\begin{itemize}
    \item 即便已经补上 shared-patch workflow 和端到端 \texttt{Conv1d} 自编码器，它们仍然只是极小的教学版表示学习系统；
    \item 模型很浅，瓶颈也极小，还不能表达更复杂的非线性流形；
    \item 数据规模、噪声模型和光谱多样性都远低于真实巡天；
    \item 我们目前虽然已经把 latent 表示接到一个最小最近邻检索与 triage 摘要，但离真实的大规模表示检索、聚类索引和跨巡天迁移还很远。
\end{itemize}

但它已经足够说明一个方向：现代 AI 的很多能力，并不来自直接预测标签，而来自先学会一个对数据结构有压缩力的内部表示。

\section{和第 23 章、第 24 章、第 32 章的关系}

可以把这几章连起来看：

\begin{itemize}
    \item 第 23 章用 PCA 和聚类说明了无监督结构探索的基本入口；
    \item 第 24 章说明了异常检测不只是“找离群点”，而是找偏离正常结构的样本；
    \item 第 32 章说明了光谱中的局部模式适合卷积表示；
    \item 本章则把这些思想汇合到“低维表示 + 重构误差”这条主线里。
\end{itemize}

\begin{warnbox}[自编码器入门中的常见误区]
\begin{itemize}
    \item 以为重构误差低就等于模型已经学会了科学机制；
    \item 忘记做 baseline，导致不知道 latent 表示到底补了什么；
    \item 在训练集中混入大量异常样本，却仍然期待重构误差自动分离异常；
    \item 只看瓶颈维度的数字，不回头检查 latent 是否真的保留了有解释力的结构；
    \item 把一个教学用小 autoencoder 直接误读成真实基础模型。
\end{itemize}
\end{warnbox}

\section{AI 助手如何帮助你}

在这个案例里，AI 助手适合帮助你：

\begin{itemize}
    \item 检查编码器、解码器和梯度更新的维度是否一致；
    \item 帮你把训练损失、测试误差和异常阈值整理成清楚的对照输出；
    \item 帮你把 latent 坐标与连续谱斜率、吸收深度之间的关系解释成更容易讲授的文字；
    \item 帮你把“为什么这是自监督入口而不是分类任务”总结成一段课程说明。
\end{itemize}

但是否真的应该把异常解释为新物理、坏数据、宇宙线污染还是标注流程问题，仍然需要你自己判断。

\section{练习}

\begin{exercisebox}[基础练习]
把 notebook 里的瓶颈维度从 2 改成 1，再改成 3，比较正常样本与异常样本的重构误差会怎样变化。解释瓶颈过小和过大分别会带来什么教学后果。
\end{exercisebox}

\begin{exercisebox}[进阶练习]
把训练集中的一部分异常样本故意混进去，再重新训练自编码器。比较异常阈值和测试样本重构误差的变化，说明“正常流形被污染”会怎样影响异常检测。
\end{exercisebox}

\begin{exercisebox}[开放练习]
结合第 32 章的一维卷积章节，设计一个后续实验方案：先用卷积层提取局部光谱特征，再把这些特征送进一个小型自编码器。说明这样做可能比直接在原始 flux bin 上训练更适合哪些任务。
\end{exercisebox}

\section{本章小结}

本章最重要的价值，不在于把学生立刻带进大规模自监督模型，而在于建立一种现代表示学习直觉：模型可以先通过重构任务学会低维 latent 表示，再利用这种表示去理解正常结构、压缩数据和发现异常。对天文和物理数据分析来说，这条主线和监督学习、卷积表示一样，都是进入现代 AI 工作流不可缺少的一部分。
